% Strong-Kleene construction from the original 1975 article.
\chapter{Partial Truth and Fixed Points}
\label{sec:p11-kripke-fixed-points}

A language can contain a predicate intended to express its own truth while
some sentences remain without a truth value. The construction below fixes
the ordinary interpretation first, assigns partial extensions to the truth
predicate, and repeatedly corrects those extensions by evaluating the
language. Its central order records the addition of information. This gives
a least fixed point, several larger fixed points, and precise distinctions
between grounded, ungrounded and paradoxical sentences
\cite[pp.~700--709]{Kripke1975}.

\section{Language, coding and partial interpretation}

\begin{hbconvention}{def:p11-ktf-base}{Base structure and codes}
Work in ordinary classical set theory. Let \(L_0\) be a countable first-order
language with equality, and let \(\mathcal M\) be a fixed classical
\(L_0\)-structure with nonempty set domain \(D\). Functions in
\(\mathcal M\) are total; its predicates are bivalent. Form the language
\(L=L_0\cup\{T\}\), where \(T\) is a new unary predicate. Its formulas
are finite strings with primitive connectives \(\neg,\land,\lor\) and
quantifiers \(\forall,\exists\). Put
\(A\to B:=\neg A\lor B\) and
\(A\leftrightarrow B:=(A\to B)\land(B\to A)\).

Fix an injective coding \(c:\operatorname{Sent}(L)\to D\) of the set of
closed formulas, write \(\ulcorner A\urcorner=c(A)\), and set
\(C=c[\operatorname{Sent}(L)]\). Fix a listing of the countable set \(C\).
A code is a domain element; the expression
\(T(\ulcorner A\urcorner)\) is object-language notation only when a
closed term naming that code has been supplied. In general, truth at a code
means the value of \(T(x)\) under an assignment sending \(x\) to that
element. The explicit arithmetic instance below supplies closed names.
No assumption that \(D\) itself is countable is made.
\end{hbconvention}

\begin{hbdefinition}{def:p11-ktf-states}{Partial truth states and information order}
A partial truth state is a pair \(S=(S_1,S_2)\) of disjoint subsets of
\(D\). Its positive extension is \(S_1\), its negative extension is
\(S_2\), and elements outside their union have no assigned truth value.
Let \(\mathcal P\) be the set of these states. Define
\[
 S\preceq R\quad\Longleftrightarrow\quad
 S_1\subseteq R_1\ \text{and}\ S_2\subseteq R_2,
 \qquad \bot_T=(\varnothing,\varnothing).
\]
Two states are compatible when
\(S_1\cap R_2=\varnothing=S_2\cap R_1\).
A chain is a subset of \(\mathcal P\) in which every two states are
comparable under \(\preceq\).

Use \(1,0,\mathbf u\) for true, false and undefined. Their information
order is \(\mathbf u\leq_i1\) and \(\mathbf u\leq_i0\), together
with reflexivity; \(0\) and \(1\) are incomparable. Here
\(\mathbf u\) records a gap in an interpretation, rather than a numerical
degree of truth. Only \(T\) varies with \(S\). In particular,
\(\bot_T\) leaves the ordinary predicates of \(\mathcal M\) fixed.
\index{information order}\index{truth-value gap}
\end{hbdefinition}

\begin{hbdefinition}{def:p11-ktf-evaluation}{Strong-Kleene evaluation}
An assignment \(g\) maps variables to \(D\). Evaluate terms in
\(\mathcal M\); write \(t^{\mathcal M,g}\) for the resulting element.
The value \(v_S(A,g)\) is defined by recursion on formulas. Equality and
\(L_0\)-atoms receive their ordinary classical values. For a truth atom,
\[
 v_S(T(t),g)=
 \begin{cases}
  1,&t^{\mathcal M,g}\in S_1,\\
  0,&t^{\mathcal M,g}\in S_2,\\
  \mathbf u,&t^{\mathcal M,g}\notin S_1\cup S_2.
 \end{cases}
\]
Negation exchanges \(1\) and \(0\), and leaves \(\mathbf u\)
undefined. The binary operations are the following complete tables.
\begin{center}
\begin{tabular}{c|ccc@{\qquad\qquad}c|ccc}
\(\land\)&\(1\)&\(0\)&\(\mathbf u\)&
\(\lor\)&\(1\)&\(0\)&\(\mathbf u\)\\\hline
\(1\)&\(1\)&\(0\)&\(\mathbf u\)&
\(1\)&\(1\)&\(1\)&\(1\)\\
\(0\)&\(0\)&\(0\)&\(0\)&
\(0\)&\(1\)&\(0\)&\(\mathbf u\)\\
\(\mathbf u\)&\(\mathbf u\)&\(0\)&\(\mathbf u\)&
\(\mathbf u\)&\(1\)&\(\mathbf u\)&\(\mathbf u\)
\end{tabular}
\end{center}
Here a row gives the left argument and a column the right argument. Write
\(g[x:=d]\) for the assignment agreeing with \(g\) except at \(x\).
Then
\[
 v_S(\forall x A,g)=
 \begin{cases}
  1,&v_S(A,g[x:=d])=1\text{ for every }d\in D,\\
  0,&v_S(A,g[x:=d])=0\text{ for some }d\in D,\\
  \mathbf u,&\text{otherwise},
 \end{cases}
\]
\[
 v_S(\exists x A,g)=
 \begin{cases}
  1,&v_S(A,g[x:=d])=1\text{ for some }d\in D,\\
  0,&v_S(A,g[x:=d])=0\text{ for every }d\in D,\\
  \mathbf u,&\text{otherwise}.
 \end{cases}
\]
The value of a sentence is independent of \(g\), and is written
\(v_S(A)\). These clauses permit a true disjunction to have an undefined
disjunct, and a false conjunction to have an undefined conjunct.
\index{strong Kleene scheme}
\end{hbdefinition}

\begin{hblemma}{thm:p11-ktf-evaluation}{proved=annex:p11.ktf.evaluation}{Evaluation and persistence}
The clauses define a unique value for every formula and assignment. This
value depends only on the assignment to the formula's free variables.
Capture-free substitution of a term for a variable agrees with changing
the assignment to the denotation of that term.
If \(S\preceq R\), then
\[
 v_S(A,g)\leq_i v_R(A,g)
\]
for every formula and assignment. In particular, a value already equal to
\(1\) or \(0\) persists when information is added.
\end{hblemma}

\section{The truth jump and the ordered space}

\begin{hbdefinition}{def:p11-ktf-jump}{Truth jump and fixed point}
The Strong-Kleene truth jump \(\Phi:\mathcal P\to\mathcal P\) is
\[
 \begin{split}
 \Phi(S)_1&=\{\ulcorner A\urcorner:A\in\operatorname{Sent}(L),\ v_S(A)=1\},\\
 \Phi(S)_2&=(D\setminus C)\cup
 \{\ulcorner A\urcorner:A\in\operatorname{Sent}(L),\ v_S(A)=0\}.
 \end{split}
\]
A fixed point is a state \(F\) satisfying \(\Phi(F)=F\).
An inflationary seed is a state \(S\) satisfying \(S\preceq\Phi(S)\).
This is an additional condition on a seed, not a property of every state.

At \(\bot_T\), even \(T\) applied to a non-sentence code is undefined.
After one jump every element of \(D\setminus C\) belongs to the negative
extension. This convention follows the original construction
\cite[pp.~702--703]{Kripke1975}.
\index{truth jump}\index{fixed point!truth}
\end{hbdefinition}

\begin{hblemma}{thm:p11-ktf-order}{proved=annex:p11.ktf.order}{Chain unions and a monotone jump}
The relation \(\preceq\) is a partial order with least element
\(\bot_T\). A family of pairwise compatible states has a least upper
bound given by the coordinatewise union. Every chain is such a family.
The map \(\Phi\) has disjoint output coordinates and is monotone:
\[
 S\preceq R\quad\Longrightarrow\quad\Phi(S)\preceq\Phi(R).
\]
The full ordered space \(\mathcal P\) is not a lattice.
\end{hblemma}

For example, assigning the same domain element positively in one state and
negatively in another leaves no consistent common extension. Therefore a
complete-lattice fixed-point theorem does not apply directly to
\(\mathcal P\). The increasing construction uses compatible unions.

\section{Transfinite construction and the closure bound}

\begin{hbconvention}{def:p11-ktf-ordinals}{Ordinal stages and set-theoretic premises}
Ordinals are the transitive sets well ordered by membership, with
\(\alpha<\beta\) meaning \(\alpha\in\beta\). A successor is
\(\alpha+1=\alpha\cup\{\alpha\}\); a nonzero ordinal with no immediate
predecessor is a limit. The first infinite ordinal is \(\omega\).
An ordinal is countable when it admits an injection into \(\mathbb N\).
Write \(\omega_1\) for the least ordinal with no such injection.
The construction uses Separation, Replacement and the ordinary ordinal
operations of ZF. Its recursion and countability argument are supplied in
the annex. The maximal-fixed-point result below additionally assumes
Zorn's principle: a nonempty partially ordered set in which every chain
has an upper bound has a maximal element.
\end{hbconvention}

\begin{hbdefinition}{def:p11-ktf-iteration}{Iteration from a seed}
For an inflationary seed \(S\), define the states \(S^\alpha\) by
\[
 S^0=S,\qquad S^{\alpha+1}=\Phi(S^\alpha),\qquad
 S^\lambda=
 \left(\bigcup_{\beta<\lambda}S^\beta_1,
       \bigcup_{\beta<\lambda}S^\beta_2\right)
 \quad(\lambda\text{ a nonzero limit}).
\]
A closure stage is a \(\theta\) for which
\(S^\theta=S^{\theta+1}\). The least such stage, when it exists, is the
closure ordinal of this iteration. Write \(\operatorname{Cl}(S)\) for
the fixed point at that first stage. Iteration from \(\bot_T\) is denoted
by \(K^\alpha\), and its resulting fixed point by \(K\).
\end{hbdefinition}

\begin{hbtheorem}{thm:p11-ktf-construction}{proved=annex:p11.ktf.construction}{Least fixed extension and countable closure}
For every inflationary seed, the displayed recursion is well defined,
all its states are consistent, and
\(S^\alpha\preceq S^\beta\) whenever \(\alpha\leq\beta\).
There is a closure stage \(\theta<\omega_1\).
The state \(\operatorname{Cl}(S)\) is the least fixed point extending
\(S\); subsequent stages are equal to it. Consequently \(K\) is the
least fixed point of \(\Phi\).

The countable bound concerns a truth predicate for the countably many
sentences of \(L\). It remains valid when \(D\) is uncountable.
\end{hbtheorem}

There are two separate reasons for the hypotheses. Information must persist
to make limit unions consistent. Countability then bounds the number of
strict additions of sentence codes. The first jump may add uncountably many
non-sentence elements to the negative extension; after that step only the
countably many sentence codes can change. An argument that counts all
domain elements would give the wrong bound for this construction.

\begin{hbmethod}{def:p11-ktf-use}{Applying the construction}
Specify the base structure, the syntax and its coding, and the evaluation
scheme. Verify that the proposed seed is below its jump. Evaluate one stage
using the clauses, put exactly the true sentence codes into the new positive
extension and the false sentence codes and non-codes into the new negative
extension, and use coordinate unions at limits. A stable stage gives a
fixed point. The theorem guarantees eventual stabilization, but does not
turn transfinite iteration over an arbitrary structure into an effective
finite algorithm. A finite run that has not stabilized establishes only
the information already acquired.
\end{hbmethod}

\section{An explicit language with self-reference}

\begin{hbdefinition}{def:p11-ktf-arithmetic}{Arithmetic coding instance}
For the examples take \(D=\mathbb N\) with its ordinary interpretation
of \(0,s,+,\times\), where \(s(n)=n+1\), and additional unary function
symbols \(d,n,q\). The symbol \(n\) here names a function on codes;
italic \(k,m,j\) below are numerical indices. Fix distinct natural-number
token indices for this finite vocabulary, punctuation and the variables
\(x_0,x_1,\ldots\), including the still uninterpreted predicate token
\(T\). Encode a word with token indices \(a_1,\ldots,a_r\) by the
natural number whose binary expansion is
\[
 1\,1^{a_1}0\,1^{a_2}0\cdots1^{a_r}0.
\]
After the initial \(1\), each zero ends one run of ones, so the word is
uniquely recovered. The empty word has code \(1\); neither it nor \(0\)
is a sentence code. Restrict this word coding to sentences to obtain \(c\).
Let \(\bar k=s^k(0)\) be the numeral naming \(k\).

Interpret \(d\) as follows. If \(k\) codes a formula \(A(x_0)\) whose
only possible free variable is \(x_0\), let \(d(k)\) be the code of
the sentence \(A(\bar k)\) obtained by replacing its free occurrences
of \(x_0\) by \(\bar k\). On all other inputs put \(d(k)=0\).
If \(k\) codes a sentence \(A\), let \(n(k)\) be the code of
\(\neg A\), and put \(n(k)=0\) otherwise. Parsing a finite word and
substituting a closed numeral define these total functions without using
the interpretation of \(T\). The function \(q\) is fixed by the
syntactic recursion in Definition~\ref{def:p11-ktf-omega-example}.

This is a fully interpreted expansion of standard arithmetic. It names
the required syntactic operations directly; no unproved claim about their
representation by formulas of a smaller arithmetic theory is needed.
\end{hbdefinition}

\begin{hbdefinition}{def:p11-ktf-diagonal}{Liar and truth-teller sentences}
Let \(b\) be the word code of \(\neg T(d(x_0))\), and let \(e\) be
the word code of \(T(d(x_0))\). Define the actual sentences
\[
 \lambda=\neg T(d(\bar b)),\qquad
 \tau=T(d(\bar e)).
\]
By the interpretation of \(d\),
\(d(b)=\ulcorner\lambda\urcorner\) and
\(d(e)=\ulcorner\tau\urcorner\). Thus, in every partial state,
the liar \(\lambda\) has the negation of the truth value assigned by
\(T\) to its own code, and the truth-teller \(\tau\) has exactly the
value assigned by \(T\) to its own code. These are metalinguistic
identities of values derived from concrete strings and terms.
\index{liar paradox}\index{truth-teller}
\end{hbdefinition}

\begin{hbdefinition}{def:p11-ktf-grounded}{Groundedness, rank and paradoxicality}
A sentence \(A\) is grounded when \(v_K(A)\) is \(1\) or \(0\).
Its rank is then the least ordinal \(\alpha\) for which
\(v_{K^\alpha}(A)\) is defined. It is ungrounded when
\(v_K(A)=\mathbf u\). It is paradoxical, relative to this language,
structure and scheme, when \(v_F(A)=\mathbf u\) for every fixed point
\(F\). The rank measures evaluation of a sentence at a stage; entry of
its code into a truth extension is a different event.
\index{grounded sentence}\index{paradoxical sentence}\index{truth!rank}
\end{hbdefinition}

\begin{hbtheorem}{thm:p11-ktf-examples}{proved=annex:p11.ktf.examples}{Grounded values and the two self-referential cases}
Every truth-free sentence has rank \(0\). If a sentence first receives a
value at stage \(\alpha\), its code first enters the corresponding
coordinate of the bottom iteration at stage \(\alpha+1\).
The liar is paradoxical. The truth-teller is undefined throughout the
bottom iteration, but is true in one fixed point and false in another.
Consequently ungroundedness does not imply paradoxicality.

At every fixed point, \(A\) and \(T\) applied to a name of its code
have the same value. Their material biconditional has value \(1\) when
that value is defined, and \(\mathbf u\) when it is undefined.
\end{hbtheorem}

\begin{hbexample}{def:p11-ktf-local-values}{Defined compounds containing a gap}
Put \(B=(0=0)\) and \(E=(0=s(0))\) in the arithmetic instance.
At the initial state,
\[
 v_{K^0}(B)=1,\quad v_{K^0}(E)=0,\quad
 v_{K^0}(\lambda)=v_{K^0}(\tau)=\mathbf u.
\]
The complete tables give
\[
 v_{K^0}(B\lor\lambda)=1,\qquad
 v_{K^0}(E\land\lambda)=0,\qquad
 v_{K^0}(\lambda\lor\neg\lambda)=\mathbf u.
\]
The first two compounds already have rank \(0\) although they contain
the truth predicate. Groundedness follows semantic dependence under the
chosen evaluation clauses, rather than the mere presence of a truth symbol.
Also \(T(0)\) is undefined at \(K^0\) and false at \(K^1\), since
\(0\notin C\).
\end{hbexample}

\section{Why limit stages can add new information}

\begin{hbdefinition}{def:p11-ktf-omega-example}{An unbounded finite truth chain}
In the arithmetic instance, let \(A_0\) be \(0=0\). Define, by ordinary
recursion on \(k\in\mathbb N\),
\[
 q(0)=\ulcorner A_0\urcorner,\qquad
 A_{k+1}=T(\overline{q(k)}),\qquad
 q(k+1)=\ulcorner A_{k+1}\urcorner.
\]
This determines the previously reserved total function \(q\) from the
word coding. It does not evaluate any truth atom. Define
\(\sigma=\forall x_0\,T(q(x_0))\).
\end{hbdefinition}

\begin{hbtheorem}{thm:p11-ktf-limit-example}{proved=annex:p11.ktf.limit-example}{A sentence settled at the first limit}
For finite \(j,k\), \(v_{K^j}(A_k)=1\) exactly when \(k\leq j\),
and is undefined otherwise. The sentence \(\sigma\) is undefined at
every finite stage and true at \(K^\omega\). Its rank is \(\omega\),
and its code first enters the positive extension at \(K^{\omega+1}\).
In particular,
\[
 \Phi\left(\bigcup_{j<\omega}K^j\right)
 \ne \bigcup_{j<\omega}\Phi(K^j),
\]
where unions of states mean coordinatewise unions.
\end{hbtheorem}

Monotonicity therefore does not supply preservation of increasing countable
unions. At each finite stage a universal sentence can still have an
undefined instance; at the union all instances can have become true.
Stopping the construction automatically at \(\omega\) would miss the
next addition. The countable closure bound in
Theorem~\ref{thm:p11-ktf-construction} is a bound, not an identification
of the exact closure ordinal of every arithmetic language.

\section{Maximal and intrinsic fixed points}

\begin{hbdefinition}{def:p11-ktf-maximal}{Maximality and intrinsic values}
A fixed point \(F\) is maximal when no distinct fixed point extends it.
It is greatest when every fixed point is below it. A fixed point is
intrinsic when it is compatible with every fixed point. A sentence has an
intrinsic value when it is defined in some intrinsic fixed point.
These definitions are relative to the same fixed jump \(\Phi\).
In the arithmetic instance put
\(\kappa=\tau\lor(\neg\tau\land\lambda)\).
\index{fixed point!maximal}\index{fixed point!intrinsic}
\end{hbdefinition}

\begin{hbtheorem}{thm:p11-ktf-maximal}{proved=annex:p11.ktf.maximal}{Maximal extensions}
Assuming Zorn's principle, every fixed point extends to a maximal fixed
point. In the arithmetic instance every maximal fixed point decides the
truth-teller, while every fixed point leaves the liar undefined. There is
no greatest fixed point. In particular, maximality does not mean that the
interpretation is total.
The sentence \(\kappa\) is ungrounded, true at some maximal fixed
points, and undefined at others.
\end{hbtheorem}

\begin{hbtheorem}{thm:p11-ktf-intrinsic}{proved=annex:p11.ktf.intrinsic}{The lattice of intrinsic fixed points}
There is a greatest intrinsic fixed point. Intrinsic fixed points, ordered
by \(\preceq\), form a complete lattice whose least element is \(K\).
For a family of intrinsic fixed points, its join is obtained by taking the
coordinatewise union and then applying \(\operatorname{Cl}\). Its meet
is the join of its intrinsic lower bounds. A raw union need not already
be a fixed point.
\end{hbtheorem}

The full space of partial states has no upper bound for two incompatible
assignments. The lattice asserted here concerns the restricted set of
intrinsic fixed points, whose compatibility and closure have been proved.

\begin{hbdefinition}{def:p11-ktf-intrinsic-example}{A self-supporting intrinsic sentence}
Let \(a\) be the word code of the formula
\[
 T(d(x_0))\lor T(n(d(x_0))),
\]
and let \(\rho\) be its instance at \(\bar a\). Thus
\(d(a)=\ulcorner\rho\urcorner\), and
\(n(d(a))=\ulcorner\neg\rho\urcorner\). The sentence says, using its
actual code, that either it or its negation is true. For comparison put
\(\chi=\tau\lor\neg\tau\).
\end{hbdefinition}

\begin{hbtheorem}{thm:p11-ktf-intrinsic-example}{proved=annex:p11.ktf.intrinsic-example}{Intrinsic truth beyond groundedness}
The sentence \(\rho\) is ungrounded but has the intrinsic value true.
The sentence \(\chi\) is true at every fixed point in which it is
defined, and false at none, but has no intrinsic value. Thus being never
false among fixed points does not suffice for having an intrinsic value.
\end{hbtheorem}

\section{Evaluation schemes and the metalanguage boundary}

\begin{hbdefinition}{def:p11-ktf-supervaluation}{Unrestricted supervaluation at a partial state}
A total extension of \(S\) is a state \(R=(R_1,R_2)\) with
\(S\preceq R\) and \(R_1\cup R_2=D\). Evaluate \(L\) classically in
each such extension. A sentence is supertrue at \(S\) when it is true in
every total extension, superfalse when false in every total extension,
and otherwise undefined. A total extension here is an arbitrary bivalent
interpretation of \(T\) extending the two coordinates. It need not be a
fixed point or the positive extension of a classically consistent theory.
\index{supervaluation}
\end{hbdefinition}

\begin{hbtheorem}{thm:p11-ktf-boundaries}{proved=annex:p11.ktf.boundaries}{Three diagnostic boundaries}
Every partial state has a total extension. The sentence
\(\lambda\lor\neg\lambda\) is supertrue at every partial state,
although it is undefined at every Strong-Kleene fixed point.

If one repeatedly revises a total interpretation by evaluating sentences
classically and assigning their resulting values to their codes, the
membership bit of the liar code flips at every successor step.
This successor revision is not an increasing information sequence.

At every Strong-Kleene fixed point, the metalinguistic assertion
\(\ulcorner\lambda\urcorner\notin F_1\) is true, while the
object-language sentence \(\neg T(d(\bar b))\) remains undefined.
\end{hbtheorem}

The original comparison also considers restrictions on the admitted total
extensions and other evaluation schemes
\cite[pp.~711--712]{Kripke1975}. For example, weak Kleene evaluation makes a
binary compound undefined whenever one argument is undefined. Accordingly
\((0=0)\lor\lambda\) is true by the Strong-Kleene table but undefined
by that weak rule. These choices alter the jump and its fixed points.

Revision from a total hypothesis uses a different process. The successor
calculation above specifies its oscillation, but a transfinite revision
theory also needs its own limit policy before stable or nearly stable
behaviour can be discussed. Likewise the supervaluation definition
quantifies over total extensions of a partial state; restricting it to
total fixed points would give an empty family in the presence of the liar.

Finally, fixed-point transparency is equality of semantic values. It does
not make a material truth biconditional true when both sides are undefined.
The metalinguistic description of a gap and the object-language negation of
truth remain distinct. This is the local form of the hierarchy problem
that persists after a partial truth predicate has been constructed
\cite[pp.~713--715]{Kripke1975}.

\section{Empirical dependence and finite language levels}

The original motivation concerns ordinary truth claims whose behaviour
depends on which sentences people have uttered. A syntactically legitimate
sentence can become paradoxical under one interpretation of an ordinary
predicate and grounded under another. Kripke does not infer that all
self-reference is illegitimate, or that a syntactic ban distinguishes
the ordinary cases from the dangerous ones
\cite[pp.~690--699, 706--708]{Kripke1975}.

\begin{hbdefinition}{def:p11-ktf-empirical}{A parameterized truth claim}
Add a bivalent unary base predicate \(P\) to the arithmetic instance. Let
\(H=\forall x_0(P(x_0)\to\neg T(x_0))\), and let
\(B=(0=0)\). The same strings and the same coding are used in three
base structures differing only in \(P\): its extension is empty,
\(\{\ulcorner B\urcorner\}\), or
\(\{\ulcorner H\urcorner\}\). Assigning a code to an ordinary
predicate is a specification of the base structure; it does not evaluate
the sentence denoted by that code.

For a separate mutual-dependence case add two bivalent predicates
\(P_1,P_2\) and form
\(H_1=\forall x(P_1(x)\to\neg T(x))\) and
\(H_2=\forall x(P_2(x)\to\neg T(x))\).
Interpret \(P_1\) by \(\{c(H_2)\}\) and \(P_2\) by
\(\{c(H_1),c(B)\}\). Each assertion includes the other among
the sentences it says are false, and the second additionally includes
the truth-free true sentence \(B\).
\end{hbdefinition}

\begin{hbtheorem}{thm:p11-ktf-empirical}{proved=annex:p11.ktf.empirical}{One syntax, three semantic outcomes}
With the empty extension, \(H\) is true of rank \(0\). With the
singleton extension containing \(B\), it is false of rank \(1\).
With the singleton extension containing \(H\), it is paradoxical.
Groundedness and rank are therefore relative to the interpretation of
the ordinary predicates, even with syntax and coding held fixed.
In the mutual-dependence case, \(H_2\) is false of rank \(1\)
and \(H_1\) is true of rank \(2\).
\end{hbtheorem}

\begin{hbdefinition}{def:p11-ktf-tarski-guards}{Finite Tarski languages inside a fixed point}
For the internal finite-hierarchy comparison assume that \(L_0\) is
syntactically adequate for the following predicates: each indicated
formation-and-translation code set is defined by a truth-free formula
\(A_r(x)\) of \(L_0\). This is an explicit syntactic resource,
additional to the injective metalinguistic coding used for the general
fixed-point theorem. Let \(L^{[0]}=L_0\), and put
\(C_1=c[\operatorname{Sent}(L^{[0]})]\).
For \(r\geq1\), the auxiliary language \(L^{[r]}\) has the base
vocabulary and a new unary token \(U_r\). Translate its formulas into
\(L\) by replacing each atom \(U_r(t)\) by
\(T(t)\land A_r(t)\), preserving all other constructors. Let
\(C_{r+1}\) be the codes of translated sentences of \(L^{[r]}\), and
choose the truth-free defining formula \(A_r\) for \(C_r\), as
provided by the stated syntactic adequacy. The definitions run successively:
\(C_1,A_1,C_2,A_2,\ldots\). They use finite-string formation and
translation, independently of truth values. Put
\(T_r(x)=T(x)\land A_r(x)\). Thus each \(T_r\) is a formula of
the original expanded truth language, using a truth-free guard already
expressible in its base. Adjoining new primitive guard predicates is
also possible when a weaker base lacks these resources, but that gives
a different base expansion and does not prove their original definability.
\end{hbdefinition}

\begin{hbtheorem}{thm:p11-ktf-tarski}{proved=annex:p11.ktf.tarski}{Finite hierarchy embedding}
At any Strong-Kleene fixed point, every \(T_r\) is bivalent on all of
\(D\), and its positive extension is exactly the set of true translated
sentences with codes in \(C_r\). Each \(L^{[r]}\), with \(U_r\)
interpreted by \(T_r\), is consequently a classical interpreted
language whose truth predicate is represented by \(T_{r+1}\).
These interpretations do not depend on which fixed point is chosen.
Every translated sentence at any finite level is grounded.
\end{hbtheorem}

This is the noncumulative finite hierarchy: each new level has the base
vocabulary and the immediately preceding truth predicate. A cumulative
hierarchy retains all preceding truth tokens instead. The original article
compares the two choices at each finite level. Their effective arithmetic
translations are supplied in Theorem~\ref{thm:p11-ktf-finite-equivalence}.
It also discusses possible extensions to transfinite levels, without
giving those transfinite definitions \cite[p.~710 and n.~26]{Kripke1975}.
The finite embedding above does not supply a transfinite hierarchy or
establish its converse characterization of sentences with an intrinsic level.
The separately specified recursive-transfinite convention and its converse
are given below in Definition~\ref{def:p11-ktf-completed-hierarchy} and
Theorem~\ref{thm:p11-ktf-uniform-converse}; their scope is not inferred
from this finite embedding.

\section{Satisfaction and signed closure}

\begin{hbdefinition}{def:p11-ktf-satisfaction}{A language containing its own satisfaction predicate}
For this construction suppose the base structure supplies a fixed injection
\(\iota:D^{<\omega}\to D\) and its decoding, and a coding
\(c_f\) of formulas of \(L_0\cup\{\mathsf{Sat}\}\) into \(D\).
Here \(\mathsf{Sat}\) is a new binary predicate. A valid argument pair
is \((\iota(s),c_f(A))\) where the finite sequence \(s\) has an
entry at every index of a free variable of \(A\).
Fix \(d_0\in D\) and extend \(s\) to an assignment \(g_s\) by
using \(d_0\) at all other indices. The value of \(A\) is independent
of those unused entries.

A satisfaction state \(Q=(Q_1,Q_2)\) consists of disjoint subsets of
\(D^2\). Evaluate ordinary atoms classically and \(\mathsf{Sat}\)
atoms by membership in these two coordinates, using the same
Strong-Kleene clauses. Its jump \(\Psi\) puts a valid argument pair
in \(\Psi(Q)_1\) or \(\Psi(Q)_2\) according as
\(v_Q(A,g_s)\) is true or false. Every invalid pair belongs to
\(\Psi(Q)_2\). The initial state has both coordinates empty.
\end{hbdefinition}

\begin{hbtheorem}{thm:p11-ktf-satisfaction}{proved=annex:p11.ktf.satisfaction}{Least satisfaction fixed point}
The jump \(\Psi\) is well defined and monotone. Its increasing ordinal
iteration from the empty state has a least fixed point \(Q^*\), satisfying
\[
 (\iota(s),c_f(A))\in Q^*_1
 \Longleftrightarrow v_{Q^*}(A,g_s)=1,
 \qquad
 (\iota(s),c_f(A))\in Q^*_2
 \Longleftrightarrow v_{Q^*}(A,g_s)=0
\]
for every valid argument pair. If \(h(X)\) denotes the least ordinal
not injectible into the set \(X\), this iteration stabilizes before
\(h(\mathcal P(D^2\times\{1,0\}))\).
The bound uses the full space of satisfaction arguments. Countability of
the formula vocabulary alone does not give the preceding \(\omega_1\)
bound when \(D\) is uncountable.
\end{hbtheorem}

The injection for sequence coding is an explicit extra resource. It was
unnecessary for a truth predicate on sentences. Kripke's original
satisfaction extension makes this distinction and notes that its closure
stage can be uncountable \cite[pp.~705--706]{Kripke1975}.

\begin{hbdefinition}{def:p11-ktf-signed}{Closure on signed sentences}
In the arithmetic instance every object has a closed numeral name. A signed
sentence is \((A,1)\) or \((A,0)\). Let \(W\) be the least set of
signed sentences satisfying the following conditions. Include the correct
sign for each closed truth-free atom, and include \((T(\bar k),0)\) when
\(k\notin C\). Close under the forcing cases of the Strong-Kleene
tables: true negation from false operand and false negation from true
operand; true conjunction from two true operands and false conjunction
from either false operand; true disjunction from either true operand and
false disjunction from two false operands. Include \((\forall x A,1)\)
from all true closed numerical instances, and its false sign from one
false instance; include \((\exists x A,1)\) from one true instance,
and its false sign from all false instances. Finally include
\((T(\overline{c(A)}),i)\) from \((A,i)\), for \(i\in\{1,0\}\).
Include replacement of a closed term in an atom by any closed term with
the same denotation in the fixed base structure. This includes replacing
a closed argument of \(T\) by its numeral; it does not identify the
two strings syntactically.
The least set exists as the intersection of all sets satisfying these
conditions; the set of all signed sentences is one such set.
\end{hbdefinition}

\begin{hbtheorem}{thm:p11-ktf-signed}{proved=annex:p11.ktf.signed}{Signed closure equals the least truth fixed point}
For the arithmetic instance,
\[
 (A,i)\in W\quad\Longleftrightarrow\quad v_K(A)=i
 \qquad(i\in\{1,0\}).
\]
In particular \(W\) has no sentence with both signs, and is also closed
under removing a truth prefix from a signed sentence when its argument is
the numeral for a sentence code.
\end{hbtheorem}

This closure presentation supplies the alternative in the original
article's note~24. It defines the same least point, while the stage
presentation additionally assigns ranks \cite[p.~706 n.~24]{Kripke1975}.

\section{Fixed points for other evaluation schemes}

\begin{hbdefinition}{def:p11-ktf-weak}{Weak Kleene and supervaluation jumps}
Weak Kleene evaluation agrees with classical evaluation on defined inputs.
A binary compound is undefined whenever either operand is undefined;
a quantified formula is undefined whenever any of its instances is
undefined. Negation preserves a gap. Write \(v^{\mathrm{wk}}_S\)
for this evaluation, and \(v^{\mathrm{sv}}_S\) for the unrestricted
supervaluation of Definition~\ref{def:p11-ktf-supervaluation}.
For each scheme form \(\Phi_{\mathrm{wk}}\) or
\(\Phi_{\mathrm{sv}}\) by the same sentence-code and non-code
clauses as \(\Phi\), using that scheme's values. Denote their
least fixed points by \(K_{\mathrm{wk}}\) and \(K_{\mathrm{sv}}\).
\end{hbdefinition}

\begin{hbtheorem}{thm:p11-ktf-schemes}{proved=annex:p11.ktf.schemes}{Monotone alternatives and comparison}
Both alternative jumps are monotone and have least fixed points reached
at countable closure stages. At each ordinal stage of their bottom
iterations, the weak Kleene state is below the Strong-Kleene state,
and the Strong-Kleene state is below the supervaluation state.
The same inclusions hold between the least fixed points.
Every classically valid sentence is supertrue at the initial state and
belongs to the positive truth extension after the first supervaluation
jump. If the base predicate \(\mathsf{PQT}\) recognizes the codes of
sentences provable in classical first-order logic, then
\(\forall x(\mathsf{PQT}(x)\to T(x))\) is supertrue at stage
\(1\), and its code enters the positive extension at stage \(2\).
Neither of the two alternative constructions makes the liar defined
at a fixed point.
\end{hbtheorem}

Groundedness and intrinsic values must always be indexed by the chosen
scheme. For example \(\tau\lor\neg\tau\) is grounded and true
under unrestricted supervaluation, whereas
Theorem~\ref{thm:p11-ktf-intrinsic-example} shows it has no intrinsic
value under Strong Kleene. None of these inclusions identifies the
schemes or establishes that one is the ordinary-language interpretation
of truth \cite[pp.~711--712]{Kripke1975}.

\begin{hbdefinition}{def:p11-ktf-restricted}{Restricted supervaluation families}
A negation-consistent completion of a state is a total extension whose
positive coordinate contains no sentence code together with the code of
its negation. This is the first consistency restriction in the original
article; it does not demand consistency of every finite set of positive
sentences. Non-sentence elements are unrestricted by this condition.
A maximal-consistent completion is a total extension whose positive
coordinate is the code set of a deductively closed, classically consistent
theory containing exactly one of each sentence and its negation.
Its positive coordinate contains no non-sentence elements. For either family,
restricted supertruth and superfalsity mean truth and falsity in all
members of that family. A state is eligible only when the indicated
family is nonempty. Universal quantification over an empty family is not
used to assign both values to a sentence.

The arithmetic function \(n\) returns the code of the negation of a
sentence, and returns \(0\) on other inputs. A bivalent base predicate
\(\mathsf{Sent}\) recognizes \(C\).
\end{hbdefinition}

\begin{hbtheorem}{thm:p11-ktf-restricted}{proved=annex:p11.ktf.restricted}{Consistency-restricted fixed points}
On their eligible states, the two restricted jumps are monotone, preserve
eligibility, and have least fixed points obtained by increasing iteration
from the empty state, with countable closure bounds.
The guarded consistency sentence
\[
 \forall x\,(\mathsf{Sent}(x)\to\neg(T(x)\land T(n(x))))
\]
is supertrue at the initial state for each restricted scheme.
The unguarded sentence \(\forall x\,\neg(T(x)\land T(n(x)))\)
is supertrue at stage \(1\) in the negation-consistent scheme, and
at stage \(0\) in the maximal-consistent scheme.
The maximal-consistent scheme also makes
\[
 \forall x\,(\mathsf{Sent}(x)\to(T(x)\lor T(n(x))))
\]
supertrue there. A sentence's code enters the positive extension at the
jump following the stated evaluation stage.
These are statements about the restricted completions;
they do not assert that some fixed point decides every sentence.
\end{hbtheorem}

The logical proof system for maximal consistency is classical, with finite
proofs; its countable sentence-extension construction is supplied in
Annex~\ref{annex:p11.ktf.restricted}.
The unguarded consistency formula is Kripke's formula~(14).
After the first jump, non-sentence inputs belong to the negative
coordinate, so all extending completions make \(T(x)\) false there.
This explains its unguarded universal form at the least fixed point
without imposing a value on the initial truth state.
The second formula is true because each admitted completion decides
every sentence; the supervaluation itself can still leave its individual
disjuncts undefined \cite[pp.~711--712]{Kripke1975}.

\section{Closing off and the residual hierarchy}

\begin{hbdefinition}{def:p11-ktf-closed}{Closing off a partial truth predicate}
Given a Strong-Kleene fixed point \(F\), put
\(\widehat F=(F_1,D\setminus F_1)\), and evaluate the language
classically with that total interpretation of \(T\). Write
\(\widehat v_F\) for its classical evaluation. Thus \(T\) now
means membership in the positive extension of the former partial model.
\end{hbdefinition}

\begin{hbtheorem}{thm:p11-ktf-closed}{proved=annex:p11.ktf.closed}{Conditional transparency after closing off}
Every sentence defined at \(F\) retains its value at \(\widehat F\).
For a name \(k\) of any sentence \(A\), the classical conditional
\[
 (T(k)\lor T(n(k)))\to(A\leftrightarrow T(k))
\]
is true at \(\widehat F\). The closed-off model is not a fixed point:
the liar is true there although its code does not belong to \(F_1\).
Its classical truth set therefore differs from \(F_1\).
\end{hbtheorem}

Closing off retains the old defined values and makes it possible to deny
that the liar belongs to the old positive truth extension. It also changes
the language's truth conditions. A new Tarskian truth predicate for this
classical language is needed to express its actual truth set; the old
predicate does not become that predicate. This is the final comparison in
the original article \cite[pp.~714--715]{Kripke1975}.

\section{Effective bounds and definability}

The original article also indicates applications to modal languages,
possible-world-dependent truth, predicates of necessity and propositional
attitudes, and impredicative substitutional quantification. It distinguishes
an operator such as necessity from a predicate applied to sentence codes,
and gives \(\Box T(x)\) and \(T(\mathsf{nec}(x))\) as different
possible definitions of a necessity predicate. The former evaluates the
truth atom across accessible worlds; the latter evaluates the code of a
modalized sentence. Those sketches do not specify a complete new jump or
prove that the definitions coincide in all modal models
\cite[pp.~712--714 and nn.~31--33]{Kripke1975}.

The preceding construction concerns arbitrary countable syntax and a fixed
set-sized base structure. The original article ends with stronger effective
claims for specifically arithmetic and acceptable structures. They use
resources that a mere countable vocabulary does not supply.

\begin{hbdefinition}{def:p11-ktf-effective}{Effective descriptive classifications}
A recursive ordinal is the order type of a computable well-ordering of a
computable subset of \(\mathbb N\). The first nonrecursive ordinal,
the supremum of these order types, is \(\omega_1^{\mathrm{CK}}\).
A lightface \(\Pi^1_1\) subset of \(\mathbb N\) is defined by a
universal set quantifier followed by an arithmetical formula, with no set
parameters. A lightface \(\Delta^1_1\) set and its complement are
both lightface \(\Pi^1_1\). A hyperarithmetical set belongs to
the effective iterated-jump class \(\mathsf{HYP}\) defined in
Method~\ref{def:p04-kleene-ordinal-notations}; the equality of that
class with \(\Delta^1_1\) is a characterization theorem, rather
than an identification assumed in the proofs below.
These classifications concern effective definability, independently of
the countable upper bound \(\omega_1\) for arbitrary syntax.
\end{hbdefinition}

\begin{hbconvention}{def:p11-ktf-effective-coding}{Effective arithmetic coding}
For the effective results choose computable token indices in the arithmetic
instance, so words, formulas and sentences are decidable sets of codes and
substitution is computable. The fixed functions \(d,n,q\) have their
specified computable interpretations. The pure language
\(0,s,+,\times,=\) also suffices: the graph lemma below supplies
formulas for the needed numerical operations. The primitive
interpretations in this vocabulary are computable. Additional
set predicates would be explicit oracle parameters in an effective
classification. All classifications here are
lightface, with no set parameters, in classical set theory.
\end{hbconvention}

\begin{hblemma}{thm:p11-ktf-arithmetic-graphs}{proved=annex:p11.ktf.arithmetic-graphs}{Arithmetic formulas for computable numerical graphs}
Every partial computable numerical function has its graph defined by an
arithmetical formula, obtained effectively from a machine for the function.
Every computable predicate, finite-word operation and substitution function
used below therefore has a truth-free arithmetic definition. For a total
function its graph has exactly one output at each standard-model input.
\end{hblemma}

\begin{hbdefinition}{def:p11-ktf-recursive-trees}{Recursive trees and heights}
Code finite sequences effectively, with computable length, entries and
concatenation. A tree is a prefix-closed subset of
\(\mathbb N^{<\omega}\) containing the empty sequence. It is recursive
when membership is computable. A branch is a function with all its finite
initial sequences in the tree; a tree with no branch is well founded.
A recursive family \(\mathcal T_n\) has membership uniformly computable
in \(n\) and the sequence code. A node with well-founded subtree has height
\[
 h_{\mathcal T}(s)=\sup\{h_{\mathcal T}(s\mathbin{{}^\frown}k)+1:
                         s\mathbin{{}^\frown}k\in\mathcal T\},
\]
defined by well-founded recursion, with the empty supremum equal to \(0\).
\end{hbdefinition}

\begin{hblemma}{thm:p11-ktf-tree-normal-form}{proved=annex:p11.ktf.tree-normal-form}{Well-founded-tree normal form and arithmetic choice}
Every lightface \(\Pi^1_1\) set \(A\) has a recursive family
\(\mathcal T_n\), effectively obtained from its definition, such that
\(n\in A\) exactly when \(\mathcal T_n\) is well founded.
An existential set assertion with an arithmetical matrix, possibly relative
to a fixed set parameter \(Z\), is equivalent to the existence of a
branch of a tree recursive relative to \(Z\). If an arithmetical relation
\(\theta(j,Y,Z)\) has a set witness for every natural \(j\), one set
can code witnesses for all \(j\). The latter is a uniform arithmetical
choice statement, not countable choice for arbitrary families of sets.
\end{hblemma}

\begin{hbtheorem}{thm:p11-ktf-pi-upper}{proved=annex:p11.ktf.pi-upper}{Upper definability bound for arithmetic truth}
The positive and negative sentence-code extensions of the least
Strong-Kleene point in the effective arithmetic instance are lightface
\(\Pi^1_1\). The same holds for unrestricted supervaluation over
total extensions of both coordinates. Thus the true and false extensions
of any fixed monadic formula are lightface \(\Pi^1_1\) for either scheme.
\end{hbtheorem}

The original two-coordinate supervaluation family is called
\(V_{\mathrm{VB}}\) in Speranski's later notation, distinguished from
his positive-coordinate-only \(V_{\mathrm{SV}}\). The proof uses our
original initial state and non-code jump
\cite[\S\S2.1,3.1]{Speranski2017TruthAuthorPreprint}.

\begin{hbdefinition}{def:p11-ktf-tree-formulas}{A diagonal family for a recursive forest}
Represent forest nodes by pairing a tree index and a finite-sequence code.
Let \(R(z,y)\) be a truth-free arithmetic definition of immediate
childhood, including the redundant conjunct \(z=z\) to retain the
node parameter syntactically. Non-node numbers have no children.
Let \(F(k,y,u)\) define
the graph of this total word operation: if \(k\) codes a formula with
only \(x_0,x_1\) possibly free, substitute \(\bar k,\bar y\)
respectively and output the sentence code; otherwise output \(0\).
After capture-free renaming of the bound variables of \(R,F\), form
\[
 B(x_0,x_1)=\forall x_2\forall x_3
 ((R(x_1,x_2)\land F(x_0,x_2,x_3))\to T(x_3)).
\]
Let \(b\) be the word code of this formula and
\(\psi_z=B(\bar b,\bar z)\). Then \(z\mapsto c(\psi_z)\)
is computable and injective, and \(F(b,y,u)\) holds exactly for
\(u=c(\psi_y)\). This defines actual strings without evaluating truth.
\end{hbdefinition}

\begin{hbtheorem}{thm:p11-ktf-tree-ranks}{proved=annex:p11.ktf.tree-ranks}{Tree height as a truth rank}
In the bottom Strong-Kleene and unrestricted-supervaluation iterations,
no \(\psi_z\) is ever false. It is defined at the least point exactly
when its subtree is well founded, and then it is true with evaluation
rank \(h_{\mathcal T}(z)\) under both schemes. Its code enters one
successor later. A node on a branch stays undefined at every stage.
\end{hbtheorem}

\begin{hbtheorem}{thm:p11-ktf-pi-representation}{proved=annex:p11.ktf.pi-representation}{Positive monadic extensions and analytical completeness}
For each of these two schemes, the positive extensions of monadic formulas
at its least arithmetic point are exactly the lightface \(\Pi^1_1\)
sets. Its positive sentence-code extension is \(\Pi^1_1\)-complete
under computable many-one reductions.
\end{hbtheorem}

\begin{hbtheorem}{thm:p11-ktf-ck-lower}{proved=annex:p11.ktf.ck-lower}{Cofinal arithmetic ranks and the Church--Kleene lower bound}
Every recursive ordinal \(\alpha\) is the evaluation rank of an actual
sentence under each of these two schemes. The closure ordinal of each
bottom iteration is at least \(\omega_1^{\mathrm{CK}}\).
\end{hbtheorem}

\begin{hbcorollary}{thm:p11-ktf-total-upper}{proved=annex:p11.ktf.total-upper}{Total monadic extensions are lightface delta-one-one}
A formula defined at every natural-number input at either least point has
a lightface \(\Pi^1_1\) positive extension and complement. Its positive
extension is consequently lightface \(\Delta^1_1\).
\end{hbcorollary}

\begin{hblemma}{thm:p11-ktf-pi-operations}{proved=annex:p11.ktf.pi-operations}{Effective positive operations on analytical definitions}
Lightface \(\Pi^1_1\) predicates are closed effectively under finite
conjunction and disjunction, number quantification and computable
substitution. These closures hold uniformly with additional number
parameters and fixed free set parameters. Consequently a positive
arithmetical formula in a set predicate, and a formula with universal
set quantifiers and such a positive arithmetical matrix, preserve
\(\Pi^1_1\) definitions effectively when substituted for that predicate.
\end{hblemma}

\begin{hbdefinition}{def:p11-ktf-pi-universal}{A universal analytical index and its strict comparison}
Use separate effective word codes for definitions in second-order
arithmetic. Let \(\mathsf U(e,z)\) assert that \(e\) codes a
lightface \(\Pi^1_1\) definition with one free number variable and
that its instance at \(z\) holds. Invalid definition codes have no
instances. This universal relation is itself \(\Pi^1_1\): evaluate
the numerical matrix by the arithmetical valuation-set relation, and
universally quantify its set argument and its valuation set.

Apply the tree normal form uniformly to \(\mathsf U\), obtaining
recursive trees \(\mathcal S_{e,z}\). On a tree put a node below
its proper prefixes, and otherwise order nodes by their first unequal
entry, using the ordinary order of natural numbers. This is its
Kleene--Brouwer order. When the tree is well founded, write
\(\rho(e,z)\) for the order type of this order.
For a coded pair \(p=(e,z)\), abbreviate
\(\mathsf U(p)=\mathsf U(e,z)\),
\(\mathcal S_p=\mathcal S_{e,z}\) and \(\rho(p)=\rho(e,z)\).
Let \(\operatorname{Below}(p,q)\), for coded index pairs \(p,q\),
mean that \(\mathcal S_p\) is well founded and that there is no
order isomorphism of the Kleene--Brouwer order of \(\mathcal S_q\)
onto an initial segment, possibly the whole order, of that of
\(\mathcal S_p\). A proposed isomorphism is coded by a set of
numerical pairs and verified arithmetically.
\end{hbdefinition}

\begin{hblemma}{thm:p11-ktf-pi-comparison}{proved=annex:p11.ktf.pi-comparison}{A recursive-ordinal analytical comparison}
The relation \(\operatorname{Below}\) is \(\Pi^1_1\). For
\(q\) with \(\mathsf U(q)\), it holds exactly when
\(\mathsf U(p)\) and \(\rho(p)<\rho(q)\). For \(q\) outside
\(\mathsf U\), it holds for every \(p\) in \(\mathsf U\) and
for no other \(p\). Each \(\rho(p)\) is a recursive ordinal.
\end{hblemma}

\begin{hbtheorem}{thm:p11-ktf-effective-induction}{proved=annex:p11.ktf.effective-induction}{An effective monotone-induction bound}
Let \(\Gamma:\mathcal P(\mathbb N)\to\mathcal P(\mathbb N)\)
be monotone. Suppose it has an effective uniform procedure converting
any \(\Pi^1_1\) definition of \(P(u,\vec z)\) into a
\(\Pi^1_1\) definition of
\(n\in\Gamma(\{u:P(u,\vec z)\})\), with the same numerical
parameters. Then its bottom iteration stabilizes by
\(\omega_1^{\mathrm{CK}}\). Every member of its least fixed
point enters at a stage strictly below that ordinal.
\end{hbtheorem}

The comparison-index argument is a numerical form of the normed-induction
construction in Moschovakis's later exposition
\cite[7C.8]{Moschovakis2009DSTAuthor}. Its uniform-definition hypothesis
is verified for the concrete operators in the next proof; it is not
inferred merely from monotonicity.

\begin{hbtheorem}{thm:p11-ktf-ck-exact}{proved=annex:p11.ktf.ck-exact}{Exact arithmetic closure ordinal}
For the effective arithmetic instance, the bottom Strong-Kleene
iteration and the unrestricted two-coordinate supervaluation iteration
both have closure ordinal exactly \(\omega_1^{\mathrm{CK}}\).
\end{hbtheorem}

\begin{hbdefinition}{def:p11-ktf-product-comparison}{An alternating comparison of two trees}
For two recursive trees in one uniform family, code a pair of nodes
\((s,t)\) with its family index by one natural number \(z\).
Let \(P(z,a)\) and \(Q(z,b)\) recognize immediate children of its
first and second nodes. Let \(J(z,a,b,w)\) be the arithmetic graph
coding the corresponding pair of children. With the same diagonal
substitution graph \(F\) as above, form
\[
 D(x_0,x_1)=\forall a\bigl(P(x_1,a)\to
   \exists b\exists w\exists u\,
   (Q(x_1,b)\land J(x_1,a,b,w)\land F(x_0,w,u)\land T(u))\bigr),
\]
with fresh bound variable names and the redundant conjunct \(z=z\)
included in \(P(z,a)\) to retain its node parameter in the strings.
Let \(d_D\) be its word code and put
\(\vartheta_z=D(\overline{d_D},\bar z)\).
These actual sentences implement the alternating condition
``for every first-tree child there is a second-tree child whose pair
sentence is true.''
\end{hbdefinition}

\begin{hbtheorem}{thm:p11-ktf-total-converse}{proved=annex:p11.ktf.total-converse}{Total definition of every lightface delta-one-one set}
Every lightface \(\Delta^1_1\) subset of \(\mathbb N\) is the
positive extension of a monadic formula defined at every input at the
least Strong-Kleene point and at the least unrestricted-supervaluation
point. Thus their total monadic extensions are exactly the lightface
\(\Delta^1_1\) sets.
\end{hbtheorem}

Tree representations may be true on their sets and undefined outside.
The product comparison supplies the additional total construction.
The bridge from this pointclass to the actual iterated-jump class is
the separately constructed characterization in
Theorem~\ref{thm:p04-doh-characterization}.

\begin{hbdefinition}{def:p11-ktf-two-finite-hierarchies}{Two classical finite Tarski hierarchies}
For the arithmetic base \(L_0\), form two separate auxiliary classical
languages with effectively coded predicate tokens. Put
\(L_N^0=L_C^0=L_0\). The predicate \(\mathsf N_r\) is true of
exactly the true sentence codes of \(L_N^{r-1}\), and false on all
other inputs; put \(L_N^r=L_0\cup\{\mathsf N_r\}\).
The predicate \(\mathsf C_r\) is true of exactly the true sentence
codes of \(L_C^{r-1}\), and false on all other inputs; put
\(L_C^r=L_0\cup\{\mathsf C_1,\ldots,\mathsf C_r\}\).
Interpret each level by ordinary classical satisfaction in standard
arithmetic with its previously specified predicates. This defines each
finite level from its predecessor, and never makes its own truth
predicate a predicate of that same level. These are auxiliary Tarski
languages; their tokens are not added to the partial truth language.
\end{hbdefinition}

\begin{hbtheorem}{thm:p11-ktf-finite-equivalence}{proved=annex:p11.ktf.finite-equivalence}{Effective equivalence of the two finite hierarchies}
For every finite \(r\), there are computable formula translations
\(u_r:L_N^r\to L_C^r\) and \(v_r:L_C^r\to L_N^r\),
preserving free variables and classical truth under every assignment
in their specified arithmetic interpretations. In particular the
languages have the same definable numerical relations at each level.
The translations use arithmetic graphs of sentence-code operations,
and provide semantic equivalence rather than syntactic inverse maps.
\end{hbtheorem}

\begin{hbtheorem}{thm:p11-ktf-effective}{proved=annex:p11.ktf.effective}{Hyperarithmetical total definability}
For the effective arithmetic instance, under Strong Kleene and under
unrestricted supervaluation over extensions of both coordinates, the
positive extensions of monadic formulas defined at every natural-number
input at the least fixed point are exactly the sets in
\(\mathsf{HYP}=\{A\subseteq\mathbb N:\exists a\in\mathcal O
\ (A\leq_T H_a)\}\), with the native notation and jump
conventions fixed in the independent arithmetic construction.
\end{hbtheorem}

These arithmetic characterizations are the results announced without
proof in the original outline \cite[pp.~715--716]{Kripke1975}.
The recursive tree, effective induction and native jump constructions
give their separate arguments. They do not identify the two schemes'
fixed points or their stagewise truth assignments.

\section{Acceptable structures and internal truth induction}

The generalization uses a different naming regime. The original
induction framework permits constants naming every domain object,
not merely a countable externally coded vocabulary
\cite[1B--1D, 5A--5B]{Moschovakis1974EIASOriginal}.
This distinction matters for an uncountable structure: its closed
instances still have internal codes, but they need not form a
countable set. The following results use that regime and its
specified elementary coding resources.

\begin{hbdefinition}{def:p11-ktf-acceptable}{An acceptable structure and its coding scheme}
Let \(\mathcal A=(A,R_1,\ldots,R_l)\) be an infinite structure
in a finite relational signature with equality. An elementary relation
means one definable by a first-order formula with finitely many
parameters from \(A\); an elementary function means a total function
with such a graph. A coding scheme consists of a subset
\(\mathsf N_{\mathcal A}\subseteq A\), a non-strict order
\(\leq_{\mathcal A}\) on that subset, and an injection
\[
 \iota:\bigcup_{n<\omega}A^n\longrightarrow A.
\]
The ordered set \((\mathsf N_{\mathcal A},\leq_{\mathcal A})\)
is isomorphic to the ordinary natural numbers with their usual order.
Write \(0_{\mathcal A},1_{\mathcal A},\ldots\) for its elements.
Let \(\mathsf{Seq}_{\mathcal A}\) be the image of \(\iota\).
For a sequence code \(s=\iota(a_1,\ldots,a_n)\), set
\(\ell_{\mathcal A}(s)=n_{\mathcal A}\) and
\(q_{\mathcal A}(s,i_{\mathcal A})=a_i\) when \(1\leq i\leq n\).
Set the length of a nonsequence to \(0_{\mathcal A}\), and set
every other projection value to \(0_{\mathcal A}\).
The structure is acceptable when \(\mathsf N_{\mathcal A}\),
\(\leq_{\mathcal A}\), \(\mathsf{Seq}_{\mathcal A}\),
\(\ell_{\mathcal A}\) and \(q_{\mathcal A}\) are elementary.
The fixed formulas defining them may themselves contain finitely
many domain parameters.

The code \(\iota()\) of the empty tuple need not equal
\(0_{\mathcal A}\), and is distinct from \(\iota(0_{\mathcal A})\).
The internal numerical markers are domain objects, not the semantic
truth values \(0,1,\mathbf u\).
\end{hbdefinition}

\begin{hbdefinition}{def:p11-ktf-general-classes}{Inductive, hyperelementary and universal-second-order relations}
For a first-order formula \(\phi(\vec x,S)\) positive in a
new \(n\)-ary relation \(S\), let its monotone operator be
\(\Gamma_\phi(B)=\{\vec a:\mathcal A\models\phi(\vec a,B)\}\).
Its least fixed point is obtained by the bottom ordinal iteration.
An inductive relation is a parameter section of such a least fixed
point: some coordinates may be fixed at finitely many specified
objects of \(A\). Denote this class by \(\operatorname{Ind}(\mathcal A)\).
Ordinary predicates may occur with either polarity; positivity is
required only of the new inductively defined predicate. A relation is
hyperelementary when it and its complement are inductive; write
\(\operatorname{HE}(\mathcal A)\) for this class.

Write \(\Pi^1_1(\mathcal A)\) for the relations definable by
a finite block of universal relation quantifiers followed by a
first-order matrix, again allowing finitely many individual parameters.
Each relation variable has a fixed finite arity and ranges over all
relations of that arity on \(A\). These are classes over the given
structure, not assertions that its domain is numerical or countable.
\end{hbdefinition}

\begin{hbconvention}{def:p11-ktf-general-language}{The name-rich language and its least truth points}
Let \(L_{\mathcal A}\) have the ordinary predicates of
\(\mathcal A\), variables \(v_1,v_2,\ldots\), and a separate
constant \(\dot a\) for every \(a\in A\), interpreted as \(a\).
Adjoin unary \(T\). Use \(\neg,\land,\exists\) as logical
primitives; \(\lor\) and \(\forall\) have their usual dual
definitions, which preserve Strong-Kleene evaluation.
Terms are variables and constants. Code them by
\(\iota(0_{\mathcal A},i_{\mathcal A})\) and
\(\iota(1_{\mathcal A},a)\), respectively. Code the remaining
constructors by
\[
\begin{array}{ll}
 R_j(t_1,\ldots,t_{n_j}):&
 \iota(2_{\mathcal A},j_{\mathcal A},c(t_1),\ldots,c(t_{n_j})),\\
 t=u:&\iota(3_{\mathcal A},c(t),c(u)),\\
 \neg B:&\iota(4_{\mathcal A},c(B)),\\
 B\land C:&\iota(5_{\mathcal A},c(B),c(C)),\\
 \exists v_i B:&\iota(6_{\mathcal A},i_{\mathcal A},c(B)),\\
 T(t):&\iota(7_{\mathcal A},c(t)).
\end{array}
\]
Here \(i\geq1\). Constant payloads are opaque domain objects;
they are not recursively parsed as codes. This is an injective code
map on the actual finite expressions.

Use the original bottom staging: \(T\) is wholly undefined on
\(A\) at stage zero; non-sentence objects first receive the false
truth-predicate value at jump one. Denote the least Strong-Kleene point
by \(K_{\mathcal A}^{\mathrm K}\), and the least unrestricted
two-coordinate-supervaluation point by \(K_{\mathcal A}^{\mathrm V}\).
Increasing compatible chain unions and the monotone jumps define
these points without a countable-code hypothesis: a strictly increasing
stage sequence cannot continue through the Hartogs ordinal of
\(\mathcal P(A\times\{0,1\})\). The countable stabilization bound
from the earlier naming regime is not asserted here.
\end{hbconvention}

\begin{hblemma}{thm:p11-ktf-internal-syntax}{proved=annex:p11.ktf.internal-syntax}{Elementary parsing and closing of formulas}
For this coding, the term, formula and sentence predicates are
elementary. So are term denotation under a coded finite assignment,
assignment update, capture-free substitution of a constant for a free
variable, and the total operation
\(\mathsf{Cl}_{\mathcal A}(e,s)\) closing a formula code \(e\)
by its coded assignment \(s\). Unspecified variables have value
\(0_{\mathcal A}\); invalid inputs to closing return a fixed
false-sentence code. Every closed formula instance is coded by this
operation. These statements assert elementary graphs, not additional
function symbols in the object language.
\end{hblemma}

\begin{hbdefinition}{def:p11-ktf-elementary-interpreter}{The positive signed elementary interpreter}
An assignment code \(s\) denotes the values
\(q_{\mathcal A}(s,i_{\mathcal A})\), with the stated default.
Use \(\mathsf V(e,s,b)\) for a formula's evaluated sign and
\(\mathsf Z(u,b)\) for a truth-predicate sign, where the domain
marker \(b=1_{\mathcal A}\) means true and \(b=0_{\mathcal A}\)
means false. The rules are these: ordinary atoms receive their ordinary
sign; a \(T\)-atom copies \(\mathsf Z\) at its term value;
negation reverses the sign; a conjunction is true if both conjuncts
are true and false if either is false; an existential is true if one
updated assignment makes its matrix true and false if all updates
make its matrix false. A sentence's \(\mathsf V\)-sign at the
empty assignment is copied into \(\mathsf Z\), and a nonsentence
object has its false \(\mathsf Z\)-sign as a base rule.
Retain all previously generated signs. All recursive occurrences of
\(\mathsf V,\mathsf Z\) in these rules are positive.

Moschovakis's classical \(\operatorname{Val}\) uses the opposite
flag convention, zero for true and one for false
\cite[5B.2]{Moschovakis1974EIASOriginal}. The displayed interpreter
uses this chapter's signs instead; its flags are not identified with
that source's flags.
\end{hbdefinition}

\begin{hblemma}{thm:p11-ktf-internal-induction}{proved=annex:p11.ktf.internal-induction}{The interpreter is one elementary positive induction}
The two signed rule relations can be packed into a single positive
elementary induction on \(A\). At its least closure,
\(\mathsf Z(u,1_{\mathcal A})\) and
\(\mathsf Z(u,0_{\mathcal A})\) are exactly the positive and
negative coordinates of \(K_{\mathcal A}^{\mathrm K}\), and
\(\mathsf V\) is exactly Strong-Kleene evaluation at that point.
Every positive or negative extension of a fixed formula is therefore
inductive over \(\mathcal A\).
\end{hblemma}

\begin{hbtheorem}{thm:p11-ktf-acceptable-ind}{proved=annex:p11.ktf.acceptable-ind}{Strong-Kleene definability over an acceptable structure}
In the name-rich language, the positive monadic extensions at
\(K_{\mathcal A}^{\mathrm K}\) are exactly
\(\operatorname{Ind}(\mathcal A)\)'s monadic relations.
The same holds in each finite arity. In particular every monadic
formula defined at every domain object has a hyperelementary positive
extension.
\end{hbtheorem}

\begin{hbdefinition}{def:p11-ktf-elementary-games}{Finite elementary-control games}
A configuration consists of a finite control label and a fixed finite
tuple of domain registers, encoded by \(\iota\). Its valid
configuration set is elementary. A finite elementary partition specifies
its type: accept, reject, deterministic successor, binary maximum,
binary minimum, domain maximum or domain minimum. Successor maps
have elementary total graphs on their indicated configuration sets;
a domain-choice map has one further argument ranging over \(A\).
Maps stay within the valid configuration set. A deterministic self-loop
is permitted.

Initially every configuration has label \(\mathbf u\).
The next label of an accept or reject configuration is \(1\)
or \(0\); deterministic transitions copy the successor's label.
Maximum and minimum configurations use their successors' labels
in the truth order \(0<\mathbf u<1\), for their two successors
or all domain choices. Compatible limit stages take the information
union. These rules are information-monotone and define a least
labeling. A configuration's resolution rank is its first stage with
label \(0\) or \(1\), if there is one; write \(\infty\)
for an unresolved configuration. Every resolution rank is a successor
ordinal, since an information union cannot introduce a new sign.
The symbol \(\infty\) here is a comparison marker above all these
ranks, not an ordinal stage.
\end{hbdefinition}

\begin{hblemma}{thm:p11-ktf-elementary-games}{proved=annex:p11.ktf.elementary-games}{From induction to games and from games to truth}
Every inductive relation is the accepted-input relation of a finite
elementary-control game which never rejects an input or a reachable
configuration. Every such game, including games with both terminal
signs, has a formula whose truth values at
\(K_{\mathcal A}^{\mathrm K}\) equal the least game labels
at the encoded initial configurations.
\end{hblemma}

\begin{hblemma}{thm:p11-ktf-selective-shuffle}{proved=annex:p11.ktf.selective-shuffle}{A totality-producing shuffle}
Let \(P\) be a game with only accepting or unresolved labels,
and \(Q\) a game with only rejecting or unresolved labels.
The alternating shuffle keeps one configuration of each and moves
the active component once before switching components. A terminal
active component returns its terminal sign. If exactly one of the
initial component configurations has a resolved label, the shuffled
configuration resolves with that sign. The shuffle is itself a finite
elementary-control game.
\end{hblemma}

\begin{hbtheorem}{thm:p11-ktf-hyperelementary}{proved=annex:p11.ktf.hyperelementary}{Total Strong-Kleene definitions are exactly hyperelementary}
The positive extensions of the monadic formulas defined at every
input at \(K_{\mathcal A}^{\mathrm K}\) are exactly
\(\operatorname{HE}(\mathcal A)\)'s monadic relations.
The same characterization holds in every finite arity and with the
specified finite individual parameters.
\end{hbtheorem}

Kripke attributes this totality result to Harrington
\cite[p.~716]{Kripke1975}. The shuffle argument supplies its
gap-resolving construction. The programme account of inductive and
hyperelementary relations gives the related alternating construction
\cite[Section~5]{HarelKozen1984IND}; neither that attribution nor
two weak Boolean representations are used as a proof of totality.

\section{General supervaluation and universal relation definitions}

The following upper-bound argument is in classical set theory with
choice. Its use is explicit: a general domain may require a family
of countermodels indexed by all its missing signed conclusions.
The arithmetic proof's canonical numerical witness selection is not
silently extended to an arbitrary domain. The coding and naming
regime remains exactly the one just specified.

\begin{hbdefinition}{def:p11-ktf-general-classical-valuation}{A classical valuation relation over the structure}
For a unary relation \(R\subseteq A\), interpret \(T\)
classically by \(R\) and its complement. Let
\(\mathsf{CV}_{\mathcal A}(R,V)\) be the first-order condition
on a ternary relation \(V\) requiring exactly one sign
\(b\in\{0_{\mathcal A},1_{\mathcal A}\}\) for every
valid formula code and finite assignment, no other tuples, and all
classical atomic and constructor evaluation equations. A truth atom
is true exactly when its term value belongs to \(R\); ordinary
atoms use their fixed predicates. Negation reverses signs, conjunction
has its classical two-sign condition, and an existential is true
exactly when some assignment update makes its matrix true.
All the syntax and update conditions in this matrix are elementary
by Lemma~\ref{thm:p11-ktf-internal-syntax}.
\end{hbdefinition}

\begin{hblemma}{thm:p11-ktf-general-classical-valuation}{proved=annex:p11.ktf.general-classical-valuation}{Uniform classical evaluation with a relation argument}
The displayed matrix is first order in its relation arguments, and
for every \(R\subseteq A\) there is exactly one ternary
\(V\) satisfying it: the actual classical formula valuation.
This is a matrix with relation arguments, not a first-order truth
predicate for the structure without such arguments.
\end{hblemma}

\begin{hbdefinition}{def:p11-ktf-general-tellers}{An independent family of internally coded truth-tellers}
Use the elementary graph
\(G(k,a,u)\) of
\(u=\mathsf{Cl}_{\mathcal A}(k,\iota(k,a))\).
Let
\[
 B(v_1,v_2)=(v_1=v_1)\land(v_2=v_2)\land
       \exists u\,(G(v_1,v_2,u)\land T(u)),
 \qquad b=c(B),\qquad \tau_a=B(\dot b,\dot a).
\]
The sentence codes \(c(\tau_a)\) are distinct, and their
graph as a function of \(a\) is elementary. The unique graph
output inside \(\tau_a\) is its own sentence code.
\end{hbdefinition}

\begin{hblemma}{thm:p11-ktf-general-tellers}{proved=annex:p11.ktf.general-tellers}{Free truth-teller coordinates at the least supervaluation point}
Every \(c(\tau_a)\) remains undefined at every stage of the
bottom unrestricted-supervaluation iteration. For any subset
\(B_0\subseteq A\), some total extension of
\(K_{\mathcal A}^{\mathrm V}\) assigns these coordinates
true exactly for \(a\in B_0\). No closure or consistency
requirement is imposed on these total truth-predicate extensions.
\end{hblemma}

\begin{hbtheorem}{thm:p11-ktf-general-sv}{proved=annex:p11.ktf.general-sv}{General supervaluation definability}
Over an acceptable structure with the specified name-rich language,
in classical set theory with choice, the positive monadic extensions
at \(K_{\mathcal A}^{\mathrm V}\) are exactly the monadic
relations in \(\Pi^1_1(\mathcal A)\). The same holds in
every finite arity with finite individual parameters. This theorem
uses unrestricted total extensions of both truth coordinates;
it does not assert the classification for the restricted-consistency
or maximal-theory completion families.
\end{hbtheorem}

The Strong-Kleene and general-supervaluation characterizations have
different right-hand classes: \(\operatorname{Ind}(\mathcal A)\)
and \(\Pi^1_1(\mathcal A)\), respectively
\cite[p.~716 and n.~36]{Kripke1975}. The arithmetic results are
their effective numerical instances, where the separate constructions
above establish the stated pointclass and hyperarithmetical equalities.
Neither classification discards the domain, naming or parameter
conditions of the general result.

\begin{hbcorollary}{thm:p11-ktf-general-finite-hierarchy}{proved=annex:p11.ktf.general-finite-hierarchy}{Elementary finite Tarski comparisons over an acceptable structure}
In the name-rich acceptable regime, the finite formation and
translation guards of Definition~\ref{def:p11-ktf-tarski-guards}
are elementary. Thus its finite fixed-point-independent Tarski
embedding applies without an additional unproved syntax-adequacy
premise. The cumulative and noncumulative finite classical Tarski
languages over \(L_{\mathcal A}\) also have value-preserving
translations with elementary total code graphs, under all domain
assignments. Invalid translation inputs can be sent to the term
code \(c(v_1)\), which is not a sentence code in any of these
languages; the numerical marker \(0_{\mathcal A}\) is not
assumed to be a noncode.
\end{hbcorollary}

\section{Recursive transfinite hierarchies and uniform level bounds}

An ordinal-indexed comparison needs both a language convention at
limits and a meaning of equivalence. Kripke explicitly leaves these
definitions outside the 1975 sketch, and reports both positive results
and technical difficulties \cite[pp.~697, 710]{Kripke1975}.
Here is one fully specified comparison. It uses recursive well-orders,
standard arithmetic with varying empirical predicates, and a new
classical truth predicate at each nonzero level, including limit levels.
It therefore does not silently identify a limit language with the
ordinary union of its predecessors.

\begin{hbconvention}{def:p11-ktf-recursive-hierarchy-base}{Fixed syntax and varying empirical interpretations}
Let \(L_{\mathrm e}\) be the language of standard arithmetic
\(0,s,+,\times\), with finitely many additional ordinary
predicates \(R_1,\ldots,R_l\) of specified arities. Let
\(\mathfrak B\) be a specified nonempty set of interpretations
of these predicates on \(\mathbb N\); the arithmetic part,
numerals, syntax and codes are held fixed throughout. The predicates
need not have computable extensions. Let
\(L_{\mathrm e,T}=L_{\mathrm e}\cup\{T\}\), with a fixed
computable finite-word coding \(c_T\); choose zero as a nonword
code, as in the earlier binary-word construction. Syntax operations
have arithmetic graphs by Theorem~\ref{thm:p04-doh-arithmetic-graphs}.
For \(\mathcal M\in\mathfrak B\), write
\(S_{\mathcal M}^\alpha\) for its original Strong-Kleene
bottom stages, retaining wholly undefined \(T\) at stage zero.
This is a separate explicit coding regime, not an alteration of the
previous examples' word-code assignments.

Fix a recursive well-order \(W=(B,\prec_W)\) of a recursive
subset of \(\mathbb N\), with \(0\in B\) its least
element. Put
\(\alpha_W(b)=\operatorname{otp}(\{d:d\prec_W b\},\prec_W)\).
The relation \(\prec_W\) is false outside \(B^2\) and
has a fixed arithmetic definition. Write \(d\preceq_W b\)
for \(d,b\in B\) and \(d=b\) or \(d\prec_W b\).
This is a linear recursive
presentation, not the branching predecessor relation on all
\(\mathcal O\)-notations. It supplies the levels represented
by this particular well-order, not all set-theoretic ordinals.
\end{hbconvention}

\begin{hbdefinition}{def:p11-ktf-completed-hierarchy}{A Tarski hierarchy completed at limit levels}
For each \(b\in B\setminus\{0\}\) introduce a distinct
auxiliary unary predicate token \(\mathsf{Tr}_b\). Set
\[
 \mathcal H_{<b}=L_{\mathrm e}\cup
       \{\mathsf{Tr}_d:0\ne d\prec_W b\},\qquad
 \mathcal H_b=L_{\mathrm e}\cup
       \{\mathsf{Tr}_d:0\ne d\preceq_W b\}.
\]
Thus \(\mathcal H_0=L_{\mathrm e}\), with no token
\(\mathsf{Tr}_0\). In the classical interpretation based on
\(\mathcal M\), \(\mathsf{Tr}_b\) is true of exactly the
codes of true \(\mathcal H_{<b}\)-sentences, and false on
every other number. Interpret these predicates by well-order recursion.
Use a separate computable word code \(c_W\) for the auxiliary
language: tagged ordinary tokens and tagged indexed-predicate tokens
are encoded into finite numerical words. A formula uses finitely
many tokens even when its language has infinitely many of them.
The sentence predicate for \(\mathcal H_{<b}\) is recursive
uniformly in \(b\), by checking its finitely many predicate
indices against \(\prec_W\).

At a nonzero limit index \(b\), the ordinary union language is
\(\mathcal H_{<b}\). The completed language \(\mathcal H_b\)
also contains its new truth predicate. These are different languages.
At the finite indices this convention is the cumulative finite
hierarchy, with the truth predicate for the base first occurring
at level one.
\end{hbdefinition}

\begin{hblemma}{thm:p11-ktf-recursive-hierarchy}{proved=annex:p11.ktf.recursive-hierarchy}{Existence of the completed classical hierarchy}
For every \(\mathcal M\in\mathfrak B\), the displayed
well-order recursion defines a unique simultaneous classical
interpretation of all the auxiliary predicates. Its syntax guards
and finite-word operations have uniform arithmetic graphs. Its
truth interpretations are not assumed to be arithmetically definable
without the auxiliary predicates.
\end{hblemma}

\begin{hbcorollary}{thm:p11-ktf-limit-separation}{proved=annex:p11.ktf.limit-separation}{Completion is not a harmless renaming of a union}
For every nonzero index \(b\), the full truth set of
\(\mathcal H_{<b}\) is not definable by a formula of that
same classical language in its specified interpretation. In particular
at a limit index, \(\mathcal H_b\)'s added truth predicate
is not already definable in the ordinary union
\(\mathcal H_{<b}\).
\end{hbcorollary}

\begin{hbtheorem}{thm:p11-ktf-transfinite-embedding}{proved=annex:p11.ktf.transfinite-embedding}{A uniform transfinite hierarchy embedding}
There is one computable, capture-free translation
\(\delta_W\) from the whole indexed auxiliary language into
\(L_{\mathrm e,T}\), preserving free variables. For each
\(b\in B\), every \(F\in\mathcal H_b\), every
\(\mathcal M\in\mathfrak B\) and assignment, its translated
formula is defined by stage \(S_{\mathcal M}^{\alpha_W(b)}\)
and has the classical value of \(F\) in the completed hierarchy.
That value persists at the least and every other Strong-Kleene
fixed point. The translation uses arithmetic syntax graphs and
the existing unary \(T\), not a new primitive for an ordinal
stage or a union of stage extensions.
\end{hbtheorem}

\begin{hblemma}{thm:p11-ktf-stage-compiler}{proved=annex:p11.ktf.stage-compiler}{Self-indexed finite syntax compilation}
There are uniformly computable auxiliary formulas
\(\mathsf P_b(x),\mathsf N_b(x)\in\mathcal H_b\), and
computable positive and negative word transformations
\(C_W^+(b,F),C_W^-(b,F)\in\mathcal H_b\), preserving
the free variables of a valid \(L_{\mathrm e,T}\)-formula.
For \(b=0\) both coordinate formulas are \(x\ne x\).
For \(b\ne0\) they have the form
\[
\begin{split}
 \mathsf P_b(x)&=\mathsf{Sent}_T(x)\land
   \exists d\exists y\,
   (d\prec_W\bar b\land G_e(1,d,x,y)\land\mathsf{Tr}_b(y)),\\
 \mathsf N_b(x)&=\exists d\,
   (d\prec_W\bar b\land
      (\neg\mathsf{Sent}_T(x)\lor
        \exists y\,(G_e(0,d,x,y)\land\mathsf{Tr}_b(y)))).
\end{split}
\]
Here \(G_e\) is the arithmetic graph of the actual total
word compiler with index \(e\), and its outputs on a valid
input formula are the words \(C_W^+,C_W^-\).
The transformations replace a truth atom by \(\mathsf P_b\)
or \(\mathsf N_b\), ordinary atoms by their positive or
negative classical condition, and the remaining constructors by
their Strong-Kleene true/false conditions. The compiler and its
index are constructed without assuming that these coordinate
formulas already define a truth stage.
\end{hblemma}

\begin{hbtheorem}{thm:p11-ktf-transfinite-stages}{proved=annex:p11.ktf.transfinite-stages}{Exact coordinate representation through recursive transfinite levels}
For every \(b\in B\) and \(\mathcal M\in\mathfrak B\),
the classical extensions of \(\mathsf P_b,\mathsf N_b\)
are exactly the two coordinates of
\(S_{\mathcal M}^{\alpha_W(b)}\). For every formula \(F\)
and assignment, \(C_W^+(b,F)\) is classically true exactly
when \(F\) is true at that partial stage; \(C_W^-(b,F)\)
is classically true exactly when \(F\) is false there.
If \(F\) is undefined there, both conditions are false.
At a limit the proof represents the actual union of lower stages;
it does not replace that union by a single lower stage or assume
that the jump preserves ordinal unions.
\end{hbtheorem}

\begin{hbdefinition}{def:p11-ktf-intrinsic-profile}{Safety, intrinsic bound and semantic equivalence}
In the fixed arithmetic syntax, call a formula \(F\)
\(\mathfrak B\)-safe when it is grounded under every
\(\mathcal M\in\mathfrak B\) and every assignment to its
finitely many free variables. This is an explicit relative convention
for intrinsically grounded, not a definition supplied by the historical
outline. When \(F\) is safe let \(r_{\mathcal M}(F,g)\)
be its first defined evaluation stage, and set
\[
 \ell_{\mathfrak B}(F)=
   \sup\{r_{\mathcal M}(F,g):\mathcal M\in\mathfrak B,
                                  \ g\text{ an assignment}\}.
\]
An assignment can here be represented by a numerical tuple on the
free variables. Every individual rank is countable, by the construction
theorem. Thus the displayed supremum exists and is at most
\(\omega_1\), although it need not be countable. For a nonsafe
formula no intrinsic bound is assigned. The bound is determined by
the specified family of possible empirical interpretations, not by
selecting one of them as the actual facts. It does not assert that
the individual ranks are identical throughout the family.

Two safe original-language formulas, on a common finite list of free
variables, are \(\mathfrak B\)-equivalent, written
\(F\sim_{\mathfrak B}G\), when their least-fixed-point values
agree under every interpretation and assignment. A safe original
formula and a classical auxiliary formula are equivalent when those
same values agree with the latter's hierarchy interpretation.
The common value may depend on the interpretation and assignment;
equivalence does not require one constant truth value across models.
It is semantic equality, not a material biconditional asserted in a
language with gaps, and not equality of word codes. An intrinsic
bound in this convention is an evaluation bound, not an invariant
of semantic equivalence or a claim of minimal expressive complexity.
\end{hbdefinition}

\begin{hbdefinition}{def:p11-ktf-uniform-level}{A represented uniform level criterion}
A fixed formula \(F\) has a uniform ground-level bound \(b\)
relative to \(\mathfrak B,W\) when, for every
\(\mathcal M\in\mathfrak B\) and every assignment,
\(F\) is defined at \(S_{\mathcal M}^{\alpha_W(b)}\).
Equivalence to an auxiliary formula means equality of these values
with its classical hierarchy values, for every such interpretation
and assignment, not equality of strings or agreement only in one
actual interpretation. The signed compiler does not need to decide
whether the uniform-bound condition holds.

This is a represented bound on the intrinsic profile just defined.
It is not automatically the meaning of the undefined phrase
``intrinsic level'' in the 1975 outline, and neither safety nor this
bound is the earlier property of being an intrinsic fixed point.
Groundedness in one interpretation alone does not supply safety
or a represented recursive-ordinal bound.
\end{hbdefinition}

\begin{hbtheorem}{thm:p11-ktf-uniform-converse}{proved=annex:p11.ktf.uniform-converse}{A proved uniform-level converse}
A formula has uniform bound \(b\) exactly when it is
\(\mathfrak B\)-safe and
\(\ell_{\mathfrak B}(F)\leq\alpha_W(b)\). In fact a safe
formula is defined at every interpretation and assignment by its
supremum bound, even when that bound is \(\omega_1\).
If \(F\) has the represented uniform bound \(b\), its effective output
\(C_W^+(b,F)\in\mathcal H_b\) is equivalent to it in
the specified sense, and also agrees with it at every
Strong-Kleene fixed point over each interpretation. Conversely
every formula of \(\mathcal H_b\) has an equivalent
translated formula with uniform bound \(b\).
Thus the two level-bounded classes have the same expressive power
at each index of this completed recursive hierarchy. No optimality
of the written target formula's syntactic level is asserted.
Each recursive ordinal can be represented by choosing an appropriate
\(W,b\). A single recursive \(W\) has recursive order type
strictly below \(\omega_1^{\mathrm{CK}}\), so this construction
does not provide a level at \(\omega_1^{\mathrm{CK}}\) or above
by varying its numerical index within that same well-order.
\end{hbtheorem}

\begin{hbtheorem}{thm:p11-ktf-intrinsic-scope}{proved=annex:p11.ktf.intrinsic-scope}{Safety is not the same as a recursive intrinsic bound}
There are a single binary empirical predicate, a specified family
\(\mathfrak B_{\mathrm{wf}}\) of its interpretations, and one
fixed sentence \(\eta\) which is safe and true throughout that
family, but satisfies
\(\ell_{\mathfrak B_{\mathrm{wf}}}(\eta)=\omega_1\).
It therefore has no uniform bound represented by any recursive
well-order. The family consists of immediate-child relations of all
well-founded numerical trees; enlarging it to include a tree with
a branch destroys the safety of this sentence.
Within that restricted family the sentence is nevertheless equivalent
to \(0=0\). Failure of a recursive evaluation bound is not a
claim of semantic nonrepresentability in a finite hierarchy.

Even over all interpretations of a unary empirical predicate,
individual ranks need not be constant for a safe sentence.
Moreover a safe sentence of intrinsic bound one can be equivalent
to a truth-free sentence of bound zero. Thus neither safety nor
semantic equivalence supplies constancy or optimality of ranks.
\end{hbtheorem}

The convention supplies a proved converse at every represented
recursive level and an explicit relationship to grounding independent
of specified empirical facts. It also exhibits the exact restriction:
safety over an arbitrary empirical family does not itself supply such
a level. A comparison for nonrecursive ordinal presentations, the
general name-rich acceptable regime, or a different limit convention
would require its own coding and translation resources. The historical
outline does not select those conventions or establish their equivalence
\cite[pp.~697, 710]{Kripke1975}. Consequently this construction
is not credited as the undefined full-general intrinsic-level theorem;
the broader source comparison remains a distinct unresolved obligation.

\section{Semantic well-orders and opaque index names}

Recursiveness of the index order is not needed if its tests are ordinary
semantic guards rather than branches of the word compiler. The next
construction supplies that mechanism, also for nonnumerical index names.
It is an extension with stated resources, not the historical outline's
unspecified definition of intrinsic level \cite[p.~710]{Kripke1975}.

\begin{hbconvention}{def:p11-ktf-semantic-order-base}{A fixed coding reduct and a codable semantic order}
Fix an acceptable coding scheme on an infinite set \(A\), with the
elementary resources and name-rich finite-signature language of
Definition~\ref{def:p11-ktf-acceptable} and
Convention~\ref{def:p11-ktf-general-language}. Work in classical set
theory with choice and put \(\kappa=\lvert A\rvert\).
Supply a nonempty subset \(I\subseteq A\), a distinguished
\(o\in I\), and a strict well-order \(\prec_I\) of \(I\)
with least element \(o\), false outside \(I^2\). Write
\(d\preceq_I b\) for \(d,b\in I\) and either \(d=b\)
or \(d\prec_I b\), and put
\(\beta_I(b)=\operatorname{otp}(\{d:d\prec_I b\},\prec_I)\).
The predicates \(I,\prec_I\) must have fixed elementary
definitions in the base or be adjoined as two explicitly additional
ordinary predicates. In the latter case the resulting finite expansion,
not the unexpanded old base, is the language being compared. The
designated \(o\) and the fixed coding parameters have constant names.
No decision algorithm for \(I\) or \(\prec_I\) is assumed.

Let \(\mathfrak F\) be a specified nonempty set of interpretations
of this fixed finite ordinary vocabulary. Hold \(A\), its coding
scheme, the formulas defining that scheme, all names and their denotations,
and \(I,o,\prec_I\) fixed across this family. The coding formulas
must continue to define the held scheme in every interpretation; other
ordinary predicates may vary subject to that condition. Let
\(L_*\) denote its ordinary name-rich language and \(L_{*,T}\)
its unary partial-truth expansion, coded by the same constructor scheme.
For \(\mathcal M\in\mathfrak F\), use
\(S_{\mathcal M}^{\alpha}\) for its original bottom
Strong-Kleene stages, initially empty in both truth coordinates.
\end{hbconvention}

\begin{hbdefinition}{def:p11-ktf-semantic-order-hierarchy}{An opaque-index completed hierarchy}
For each \(b\in I\setminus\{o\}\) introduce a unary token
\(\mathsf{Tr}^{I}_b\). Define
\[
 \mathcal J_{<b}=L_*\cup
       \{\mathsf{Tr}^{I}_d:o\ne d\prec_I b\},\qquad
 \mathcal J_b=L_*\cup
       \{\mathsf{Tr}^{I}_d:o\ne d\preceq_I b\}.
\]
In each interpretation, \(\mathsf{Tr}^{I}_b\) is to be the
classical truth set of \(\mathcal J_{<b}\)-sentences and false
on all other domain objects. The new token is also supplied at a
nonzero limit index. The index order has set order type
\(\theta\), with \(\lvert\theta\rvert\leq\kappa\) and
\(\theta<\kappa^+\).

For formation first use a raw auxiliary alphabet with an indexed token
at every \(b\in A\). Its new atomic constructor is
\[
 c_I(\mathsf{Tr}^{I}_b(t))=
       \iota(8_{\mathcal A},b,c(t)).
\]
All ordinary term and formula constructors are retained; the original
\(T\)-constructor is excluded from the raw auxiliary language.
Here the index payload \(b\) is opaque, just like a constant's
payload, and is not recursively parsed. Raw formation is separated
from legality in \(\mathcal J_{<b}\) or \(\mathcal J_b\).
Those legality guards test the finitely many indexed primes using
the ordinary predicates \(I,\prec_I\). They do not require a
compiler to decide the order or replace a fixed indexed token by a
variable-indexed truth predicate in the interpreted language.
\end{hbdefinition}

\begin{hblemma}{thm:p11-ktf-semantic-order-syntax}{proved=annex:p11.ktf.semantic-order-syntax}{Elementary opaque-index syntax and prospective-template compilers}
Raw auxiliary formation, its free-variable lists, and the lower and
nonstrict sentence guards
\(\mathsf{Sent}^{I}_{<}(b,x),\mathsf{Sent}^{I}_{\leq}(b,x)\)
are elementary. There are actual total elementary syntax transformations
with graphs \(E_I(k,e,y)\) and
\(\mathsf C_I(k,\sigma,b,e,y)\), with these specifications.

The first copies an original-language prospective three-slot template
\(k\) at every raw auxiliary truth prime
\(\mathsf{Tr}^{I}_d(t)\), substituting the names of \(k,d\)
and the term \(t\) into its designated slots. The copied template
is not itself retranslated. Its admissible templates have exactly the
three designated indices free, so the argument slot is actually retained.
All other constructors are preserved.

The second produces the positive or negative Strong-Kleene condition
word for an original formula \(e\), with selector
\(\sigma=1_{\mathcal A}\) or \(0_{\mathcal A}\).
For a prospective ordinary five-slot relation-template code \(k\),
with exactly those designated indices free,
write \(K_k(k,\sigma,d,x,y)\) for its capture-free syntactic
instance with the named code \(k\) in its first slot. At a truth
prime use the following one-free-variable factory words, substituting
the prime's term for \(x\): at \(b=o\) both are \(x\ne x\);
otherwise they are
\[
\begin{split}
 p_I(k,b,x)&=\mathsf{Sent}_{*,T}(x)\land
   \exists d\exists y\,
    (I(d)\land d\prec_I\dot b\land
     K_k(k,1_{\mathcal A},d,x,y)\land\mathsf{Tr}^{I}_b(y)),\\
 n_I(k,b,x)&=\exists d\,
   (I(d)\land d\prec_I\dot b\land
    (\neg\mathsf{Sent}_{*,T}(x)\lor
       \exists y(K_k(k,0_{\mathcal A},d,x,y)
                                  \land\mathsf{Tr}^{I}_b(y)))).
\end{split}
\]
The expressions \(K_k\) denote actual copied formulas, not new
semantic relation primitives. The remaining signed transformations
use the constructor table in Lemma~\ref{thm:p11-ktf-stage-compiler}.
They form words for every domain index without deciding \(I\) or
its order. For valid templates and formulas the transformations preserve
free variables; invalid inputs have a specified false-word output.
\end{hblemma}

\begin{hbtheorem}{thm:p11-ktf-semantic-order-embedding}{proved=annex:p11.ktf.semantic-order-embedding}{Completed hierarchies over semantic well-orders}
Well-order recursion gives a unique classical interpretation of the
completed hierarchy over each \(\mathcal M\in\mathfrak F\).
At any nonzero index its added full truth set is not definable in
that same prior language \(\mathcal J_{<b}\).
There is one capture-free translation \(\Delta_I\) with a total
elementary code graph from the raw auxiliary language into
\(L_{*,T}\). For every \(b\in I\), every
\(F\in\mathcal J_b\), interpretation and assignment,
\(\Delta_I(F)\) is defined by
\(S_{\mathcal M}^{\beta_I(b)}\) and has exactly the classical
hierarchy value. That value is the same at every Strong-Kleene fixed
point over that interpretation. This holds at successor and limit
indices of arbitrary set well-orders satisfying the stated coding and
ordinary-order resources, not only recursive orders.
\end{hbtheorem}

\begin{hbtheorem}{thm:p11-ktf-semantic-order-stages}{proved=annex:p11.ktf.semantic-order-stages}{Exact coordinates from a named compiler graph}
Choose an actual ordinary formula defining the preceding elementary
graph \(\mathsf C_I\), with its five designated slots, and let
\(h\in A\) be that formula's original code. Set
\[
 P^I_b(x)=p_I(h,b,x),\qquad N^I_b(x)=n_I(h,b,x).
\]
Thus the copied relation at a nonzero index is the actual
\(\mathsf C_I(\dot h,\dot\sigma,d,x,y)\), not a relation
postulated to recognize truth stages. The compiler was constructed
before this choice of its own named formula code. It performs finite
syntax operations on that code and never interprets or recursively
recompiles the graph formula it copies.

For every \(b\in I\) and \(\mathcal M\in\mathfrak F\),
the classical extensions of \(P^I_b,N^I_b\) are exactly the
coordinates of \(S_{\mathcal M}^{\beta_I(b)}\). The actual
positive and negative outputs \(U_b^+(F),U_b^-(F)\) of
\(\mathsf C_I(h,\sigma,b,c(F),y)\) are formulas in
\(\mathcal J_b\), preserving free variables, and are classically
true exactly for the corresponding partial values of \(F\) at
that stage. Both are false when the value is undefined. No numerical
program-index theorem or primitive stage predicate is assumed.
\end{hbtheorem}

\begin{hbcorollary}{thm:p11-ktf-semantic-order-converse}{proved=annex:p11.ktf.semantic-order-converse}{Uniform converse and the actual ordinal reach}
In the held name-rich regime each bottom iteration stabilizes before
\(\kappa^+\). For a formula safe throughout \(\mathfrak F\),
define its intrinsic evaluation bound by the same supremum convention
as before, now using finite domain assignments; write
\(\ell_{\mathfrak F}(F)\). It is at most \(\kappa^+\),
and a fixed index \(b\) is a uniform bound exactly when
\(\ell_{\mathfrak F}(F)\leq\beta_I(b)\).
At every such index, \(U_b^+(F)\) is an equivalent classical
\(\mathcal J_b\)-formula, over every interpretation and
assignment; conversely every classical \(\mathcal J_b\)-formula
has a translated original-language representative with that uniform
bound. Equivalence here compares least-fixed-point and classical
values over the same held family, not codes or minimum expressive levels.

For \(A=\mathbb N\) with ordinary arithmetic coding, this
includes every countable ordinal in some supplied well-order, even
when the order is nonrecursive. The same conclusions hold in the
earlier finite-word language of standard arithmetic with its function
terms, after supplying ordinary \(I,\prec_I\) predicates for a
well-order of a numerical subset with least element zero; the raw
word transformations remain computable. More generally every ordinal
\(\alpha<\kappa^+\) can occur as an index rank in a separately
supplied order on a subset of \(A\). This requires that the
order predicates be elementary or explicitly added to the ordinary
base; it does not make arbitrary such orders definable in the
unexpanded original structure. A fixed order represents only its own
initial segments. It supplies no index of rank \(\kappa^+\)
and no comparison for a larger order that cannot be injected into
the fixed code/name domain.
\end{hbcorollary}

In the standard numerical finite-word specialization, raw parsing and
the prospective-template transformations are still computable: a predicate
index is just a numerical token payload, and the order test is copied
as an ordinary formula, not executed. Their computation graphs can thus
be arithmetical even when the supplied order relation is nonrecursive.
The language-legality guards are then formulas in the expanded ordinary
base containing that order predicate. This does not assert an arithmetic
definition of its nonrecursive extension in the unexpanded base.

If the order is varied along with the empirical interpretation, the
same prospective-template constructions still work in every structure
whose order is a well-order and whose coding scheme is the held one.
The embedding and coordinate equalities then use that structure's
\(\beta_I(b)\); a single fixed ordinal bound must not be inferred
from one shared index name. The preceding supremum criterion expressly
uses a fixed order. Neither form assigns a historical meaning to the
outline's undefined intrinsic-level term or removes the distinction
between a result in an expanded base and one in the original base.

\begin{hbtheorem}{thm:p11-ktf-expansion-transport}{proved=annex:p11.ktf.expansion-transport}{Stage-preserving recoding into an ordinary expansion}
Let an acceptable name-rich ordinary base be expanded by finitely many
ordinary predicates, keeping its domain, coding scheme, old predicate
interpretations and all names fixed. Use compatible constructor codes
\(c_0,c_1\), retaining the old predicate indices and adding new ones.
Both source and target have their own unary truth predicate for their
whole respective sentence languages. There is a capture-free translation
\(\Lambda\) from source formulas to target formulas with a total
elementary code graph such that, at every ordinal stage and assignment,
\[
 v_{S^{\alpha}_{\mathrm{new}}}(\Lambda(F),g)
       =v_{S^{\alpha}_{\mathrm{old}}}(F,g).
\]
It preserves all groundedness and first evaluation ranks, and compares
the source and target least-fixed-point values, not the values of the
unchanged source string in the enlarged target language. The ordinary
expansion's extra predicates need not be elementary in the old base.

The construction uses a prospective ordinary graph-template compiler
\(\mathsf R(k,e,y)\). If \(q_R\) is the actual target code
of that graph formula, a source truth prime \(T(t)\) is replaced by
\(\exists y(\mathsf R(\dot q_R,t,y)\land T(y))\).
The selected output on a source sentence is the code of its actual
translation; the selected output on a source nonsentence is a genuine
target nonsentence, not a false-sentence code. Thus old negative noncodes
appear at the original first-jump stage. In the numerical finite-word
version the transformation and its graph are computable and arithmetical.

Consequently an old formula with a represented uniform grounding bound
can be transported to the supplied semantic-order expansion and then
represented in its completed hierarchy at that same ordinal bound.
This does not assert that the new ordinary order predicates, or all
formulas of their hierarchy, have equivalents in the unexpanded old base.
\end{hbtheorem}
